% Generated draft: requires mathematical and source review.
\documentclass[11pt]{article}
\usepackage[margin=25mm]{geometry}
\usepackage{fontspec}
\setmainfont{Latin Modern Roman}
\usepackage{amsmath,amssymb,mathtools,unicode-math}
\setmathfont{Latin Modern Math}
\usepackage{enumitem,needspace,xurl,hyperref}
\setlength{\parindent}{0pt}
\setlength{\parskip}{5pt}
\setlength{\emergencystretch}{3em}
\newcommand{\N}{\mathbb{N}}
\newcommand{\Z}{\mathbb{Z}}
\newcommand{\Q}{\mathbb{Q}}
\newcommand{\R}{\mathbb{R}}
\newcommand{\C}{\mathbb{C}}
\newcommand{\F}{\mathbb{F}}
\newcommand{\A}{\mathbb{A}}
\newcommand{\PP}{\mathbb{P}}
\newcommand{\E}{\mathbb{E}}
\newcommand{\eps}{\varepsilon}
\newcommand{\dd}{\,\mathrm{d}}
\begin{document}
\section*{UnsolvedMath ID: 30005184}
\textbf{Source problem code:} catalogue entry
\subsection{Definitions and mathematical statement}
A compact Lie group is a finite-dimensional real Lie group whose underlying topological space is compact. For such a group \(G\), an equivariant cohomology theory is a cohomology theory incorporating a \(G\)-action; the supplied record does not specify its category of inputs, grading, coefficient objects, or axioms.

\emph{Requested construction.} Find examples of constructions of equivariant cohomology theories with good finiteness properties that apply to every compact Lie group \(G\). The registry reformulates this as requiring ``strong finiteness properties uniformly for all compact Lie groups.''

\emph{Unresolved formulation.} Neither ``good finiteness properties'' nor ``strong finiteness properties'' is defined in the supplied record. It does not identify which objects must be finite, the relevant meaning of finiteness, or any restrictions on inputs or coefficients. Moreover, ``uniformly'' could mean a common construction or common finiteness condition for every \(G\), but no bound independent of \(G\) is specified. The record also does not state whether compatibility between theories for different groups is required. Thus the universal scope over compact Lie groups is explicit, but the requested finiteness condition is not a mathematical predicate that can be stated precisely from this evidence. The question remains a request for examples, rather than a specified classification problem.

\emph{Context supplied with the question.} The original excerpt immediately continues into a section entitled ``Equivariant Elliptic Cohomology.'' It describes replacing the multiplicative group \(\mathbb G_m\) in the story of equivariant \(K\)-theory by a fixed elliptic curve \(E\), with the \(S^1\)-equivariant theory taking values in quasi-coherent sheaves on \(E\). The excerpt attributes motivations to elliptic genera, naming Miller and Rosu, and to geometric representation theory. It attributes an initial construction over \(\C\) for connected compact Lie groups to Grojnowksi, and a later coherent theory over \(\C\) for all compact Lie groups to Berwick-Evans and Tripathy. It describes Lurie's contribution as a sketch allowing coefficients beyond \(\C\), and says that the work being presented follows that outline, uses Lurie's spectral algebraic geometry, and draws inspiration from the axiomatics of Ginzburg--Kapranov--Vasserot. These are supplied contextual attributions, not independently verified claims or additional hypotheses of the question. The passage ends mid-sentence, so it cannot resolve the missing finiteness condition or establish whether elliptic cohomology is the intended exclusive setting.
\subsection{Short English statement}
Find examples of equivariant cohomology theories with suitable finiteness properties for every compact Lie group. The source does not define those properties or the intended uniformity, so a precise construction target requires further source clarification. The registry labels the problem open; this draft does not independently verify that status.
\subsection{Sources}
\begin{itemize}
\item Local registry: \texttt{data/unsolvedmath/problems.json}, permanent UnsolvedMath ID \(30005184\), source problem code \texttt{catalogue entry}, title ``Finite Equivariant Cohomology for Compact Lie Groups.'' This draft uses the supplied complete registry record. No individual proposer is identified; the supplied proposed year is \(2022\).
\item Supplied report citation: \emph{Topologie (2023)}, Oberwolfach Reports, DOI \texttt{10.4171/owr/2022/34}, pp.~1927--1986. Source URL and literature-source link labelled ``original Oberwolfach report'': \url{https://doi.org/10.4171/owr/2022/34}. The record describes the report as documenting workshop discussions and open problems. Its citation uses \(2023\), while its DOI and proposed-year metadata use \(2022\); these supplied labels are retained.
\item Supplied literature source: \emph{Proper equivariant stable homotopy theory}, \url{https://arxiv.org/abs/1908.00779}. The record does not identify a particular result from this reference that supplies the missing finiteness condition.
\item The excerpt contains internal reference labels [1] for Berwick-Evans and Tripathy, [2] for Ginzburg--Kapranov--Vasserot, [3] for Grojnowksi, and [4] for Lurie. Their corresponding bibliography entries are not supplied. It also names Miller and Rosu in discussing motivations. These attributions are retained without inventing bibliographic details or identifying any of these internal labels with the supplied arXiv link.
\item Literature assessment date: \texttt{2026-08-22}. This is supplied dated curation, not an independently verified literature assessment. It labels the problem open and describes the question as a broad request for examples, reporting examples and partial constructions in the report but no complete classification or definitive answer for all compact Lie groups. The background explicitly asks that time-sensitive status and cited claims be verified independently.
\item The registry assigns the statement status \texttt{corrected\_verified} and reports that it was reconstructed and checked against the official Oberwolfach report. This is supplied provenance; no source visit or independent verification is claimed here.
\end{itemize}
\end{document}
